\chapter{Intelligent Parameter Inference from BRENDA}
\label{ch:brenda}

\section{Introduction}
\label{sec:brenda-intro}

\textbf{Kinetic parameters} ($V_{\max}$, $K_m$, $k_{cat}$, $K_i$) determine reaction \textbf{rates and regulation} in Bio-PN models. Manually finding parameters for dozens of enzymes is:

\begin{enumerate}
    \item \textbf{Time-consuming}: Literature search takes hours per enzyme
    \item \textbf{Inconsistent}: Different studies report different values (organism, pH, temperature)
    \item \textbf{Incomplete}: Many enzymes lack published parameters
\end{enumerate}

This chapter presents \textbf{automatic parameter inference} from \textbf{BRENDA} (BRaunschweig ENzyme DAtabase), the largest enzyme kinetics repository.

\subsection{BRENDA Database}
\label{subsec:brenda-database}

\textbf{BRENDA} (\url{www.brenda-enzymes.org}) provides:

\begin{itemize}
    \item \textbf{83,000+ enzymes} classified by EC number
    \item \textbf{2.7 million} kinetic parameters ($K_m$, $k_{cat}$, $K_i$, optimal pH, etc.)
    \item \textbf{Organism-specific data}: Yeast, human, bacteria
    \item \textbf{SOAP API}: Programmatic access (requires free registration)
\end{itemize}

\subsection{SHYpn Parameter Inference}
\label{subsec:shypn-parameter-inference}

\begin{enumerate}
    \item \textbf{EC number lookup}: User specifies enzyme (e.g., EC 2.7.1.1 = hexokinase)
    \item \textbf{Organism filtering}: Prioritize yeast (\textit{S. cerevisiae}), then human, then all
    \item \textbf{Statistical aggregation}: Compute median $K_m$ (robust to outliers)
    \item \textbf{Context-aware heuristics}: Adjust by substrate, pH, temperature
    \item \textbf{Local caching}: Build SQLite database of queried parameters (offline use)
\end{enumerate}

\textbf{Benefits}:

\begin{itemize}
    \item \textbf{Speed}: Fetch all 10 glycolysis enzyme parameters in $<10$ seconds
    \item \textbf{Accuracy}: Median of 50+ literature values (robust estimate)
    \item \textbf{Reproducibility}: Cached parameters ensure consistency across sessions
\end{itemize}

\section{BRENDA Data Structure}
\label{sec:brenda-data-structure}

\subsection{Enzyme Commission (EC) Numbers}
\label{subsec:ec-numbers}

\textbf{EC classification} (hierarchical 4-level code):

\begin{itemize}
    \item \textbf{EC 2.x.x.x}: Transferases (transfer functional groups)
    \item \textbf{EC 2.7.x.x}: Transferring phosphorus-containing groups
    \item \textbf{EC 2.7.1.x}: Phosphotransferases with alcohol group as acceptor
    \item \textbf{EC 2.7.1.1}: Hexokinase (ATP:D-hexose 6-phosphotransferase)
\end{itemize}

\textbf{Examples}:

\begin{itemize}
    \item EC 2.7.1.1: Hexokinase (\ce{glucose + ATP -> G6P + ADP})
    \item EC 2.7.1.11: 6-Phosphofructokinase (\ce{F6P + ATP -> F-1,6-BP + ADP})
    \item EC 1.2.1.12: Glyceraldehyde-3-phosphate dehydrogenase (GAPDH)
\end{itemize}

\subsection{Kinetic Parameters in BRENDA}
\label{subsec:brenda-kinetic-parameters}

\textbf{Key parameter types}:

\begin{table}[h]
\centering
\begin{tabular}{@{}llll@{}}
\toprule
\textbf{Parameter} & \textbf{Definition} & \textbf{Units} & \textbf{Example} \\
\midrule
$K_m$ & Michaelis constant (substrate affinity) & mM, $\mu$M & $K_m$(glucose) = 0.1 mM \\
$k_{cat}$ & Turnover number (max reactions/enzyme/s) & $s^{-1}$ & $k_{cat}$ = 100 $s^{-1}$ \\
$K_i$ & Inhibition constant (inhibitor affinity) & mM, $\mu$M & $K_i$(ATP) = 5.0 mM \\
$V_{\max}$ & Maximum rate ($k_{cat} \cdot [E]_0$) & mM/s & $V_{\max}$ = 1.0 mM/s \\
\bottomrule
\end{tabular}
\caption{Kinetic parameter types in BRENDA}
\label{tab:brenda-parameters}
\end{table}

\textbf{Challenges}:

\begin{enumerate}
    \item \textbf{Heterogeneity}: Values vary by organism, substrate, conditions
    \item \textbf{Outliers}: Some entries are erroneous (typos, unit errors)
    \item \textbf{Incompleteness}: Not all enzymes have $K_m$ for all substrates
\end{enumerate}

\section{BRENDA SOAP API}
\label{sec:brenda-soap-api}

\subsection{Authentication}
\label{subsec:brenda-authentication}

\textbf{BRENDA API requires free credentials}:

\begin{enumerate}
    \item Register at \url{https://www.brenda-enzymes.org/}
    \item Receive username and password via email
    \item Store in SHYpn config file (\texttt{\textasciitilde/.config/shypn/brenda\_credentials.json})
\end{enumerate}

\subsection{BRENDAConnector Class}
\label{subsec:brenda-connector}

\textbf{Responsibilities}:

\begin{itemize}
    \item Send SOAP requests to BRENDA API
    \item Parse XML responses
    \item Extract numeric values (handle ranges, scientific notation, inequalities)
\end{itemize}

\textbf{Key methods}:

\begin{itemize}
    \item \texttt{get\_km\_values(ec\_number)}: Fetch all $K_m$ values for enzyme
    \item \texttt{get\_kcat\_values(ec\_number)}: Fetch all $k_{cat}$ values
    \item \texttt{get\_ki\_values(ec\_number)}: Fetch all $K_i$ values
\end{itemize}

\subsection{Example: Fetching Hexokinase $K_m$}
\label{subsec:brenda-hexokinase-example}

\begin{verbatim}
connector = BRENDAConnector(username="user@example.com", password="pass123")

km_values = connector.get_km_values("2.7.1.1")
print(f"Found {len(km_values)} Km values for hexokinase")

# Filter for yeast and glucose
yeast_glucose_km = [
    entry for entry in km_values
    if 'cerevisiae' in entry['organism'].lower() 
    and 'glucose' in entry['substrate'].lower()
]

print(f"Yeast-specific glucose Km values: {len(yeast_glucose_km)}")
for entry in yeast_glucose_km[:5]:
    print(f"  {entry['value']} {entry['unit']} ({entry['literature']})")

# Output:
# Found 347 Km values for hexokinase
# Yeast-specific glucose Km values: 23
#   0.12 mM (PubMed:12345678)
#   0.15 mM (PubMed:23456789)
#   0.08 mM (PubMed:34567890)
#   0.20 mM (PubMed:45678901)
#   0.11 mM (PubMed:56789012)
\end{verbatim}

\section{Statistical Aggregation}
\label{sec:statistical-aggregation}

\subsection{Median as Robust Estimator}
\label{subsec:median-estimator}

\textbf{Problem}: BRENDA values have \textbf{high variance} (10$\times$ difference between min/max)

\textbf{Causes}:

\begin{enumerate}
    \item \textbf{Experimental error}: Different labs, methods
    \item \textbf{Biological variation}: Different strains, growth conditions
    \item \textbf{Data entry errors}: Typos, unit confusion ($\mu$M vs. mM)
\end{enumerate}

\textbf{Solution}: Use \textbf{median} (50th percentile) instead of mean

\begin{itemize}
    \item \textbf{Median is robust to outliers} (1 extreme value doesn't skew result)
    \item \textbf{Mean is sensitive} (one 100$\times$ outlier pulls average up)
\end{itemize}

\textbf{Example} (hexokinase $K_m$ for glucose in yeast):

\begin{verbatim}
Raw values: [0.08, 0.10, 0.12, 0.15, 0.18, 0.20, 0.25, 5.0]
                                                        ↑ outlier
Mean:   0.76 mM  (pulled up by outlier)
Median: 0.165 mM (robust)
\end{verbatim}

\subsection{Confidence Intervals}
\label{subsec:confidence-intervals}

\textbf{95\% confidence interval}: Range containing true value with 95\% probability

\textbf{Bootstrap method} (non-parametric):

\begin{enumerate}
    \item Resample data with replacement (1000 times)
    \item Compute median of each resample
    \item Take 2.5th and 97.5th percentiles of resampled medians
\end{enumerate}

\textbf{Example}:

\begin{verbatim}
km_values = [entry['value'] for entry in yeast_glucose_km]
stats = calculate_statistics(km_values)

print(f"Km(glucose) for yeast hexokinase:")
print(f"  Median: {stats['median']:.3f} mM")
print(f"  95% CI: [{stats['confidence_interval_95_lower']:.3f}, "
      f"{stats['confidence_interval_95_upper']:.3f}] mM")
print(f"  Based on {stats['count']} measurements")

# Output:
# Km(glucose) for yeast hexokinase:
#   Median: 0.165 mM
#   95% CI: [0.100, 0.220] mM
#   Based on 23 measurements
\end{verbatim}

\section{Context-Aware Heuristics}
\label{sec:context-aware-heuristics}

\subsection{Organism Priority}
\label{subsec:organism-priority}

\textbf{Challenge}: BRENDA contains data from \textbf{5000+ organisms}, but user cares about specific organism (e.g., yeast)

\textbf{SHYpn strategy}: \textbf{Hierarchical filtering}:

\begin{enumerate}
    \item \textbf{Primary}: User-specified organism (\textit{S. cerevisiae})
    \item \textbf{Secondary}: Related organisms (other fungi)
    \item \textbf{Tertiary}: All organisms (if $<5$ primary values)
\end{enumerate}

\subsection{Quality Filtering}
\label{subsec:quality-filtering}

\textbf{Some BRENDA entries are low-quality}:

\begin{itemize}
    \item No literature citation
    \item Vague commentary (``approximate value'')
    \item Extreme outliers (100$\times$ different from median)
\end{itemize}

\textbf{Quality scoring} (0.0--1.0):

\begin{itemize}
    \item Has PubMed citation: +0.4
    \item Has commentary with conditions (pH, temperature): +0.2
    \item Value within reasonable range (0.01--10 mM for $K_m$): +0.4
\end{itemize}

\textbf{Filtering}: Keep only high-quality entries (score $\geq$ 0.6)

\section{Local Database Caching}
\label{sec:local-database-caching}

\subsection{Motivation}
\label{subsec:caching-motivation}

\textbf{Problem}: BRENDA API is slow (500 ms -- 2 seconds per query)

\textbf{Solution}: Cache results in local SQLite database

\begin{itemize}
    \item First query: Fetch from BRENDA + save to DB (2 seconds)
    \item Subsequent queries: Read from DB (10 ms) $\to$ \textbf{200$\times$ faster}
\end{itemize}

\subsection{Database Schema}
\label{subsec:brenda-database-schema}

\textbf{File}: \texttt{\textasciitilde/.config/shypn/brenda\_cache.db}

\textbf{Tables}:

\begin{verbatim}
CREATE TABLE brenda_raw_data (
    id INTEGER PRIMARY KEY AUTOINCREMENT,
    ec_number TEXT NOT NULL,
    parameter_type TEXT NOT NULL,  -- 'Km', 'kcat', 'Ki'
    value REAL NOT NULL,
    unit TEXT,
    substrate TEXT,
    organism TEXT,
    literature TEXT,
    commentary TEXT,
    quality_score REAL,
    query_date TIMESTAMP DEFAULT CURRENT_TIMESTAMP,
    UNIQUE(ec_number, parameter_type, value, organism, 
           substrate, literature)
);

CREATE TABLE brenda_statistics (
    id INTEGER PRIMARY KEY AUTOINCREMENT,
    ec_number TEXT NOT NULL,
    parameter_type TEXT NOT NULL,
    organism TEXT DEFAULT 'all',
    substrate TEXT DEFAULT 'all',
    median_value REAL,
    mean_value REAL,
    std_dev REAL,
    confidence_interval_95_lower REAL,
    confidence_interval_95_upper REAL,
    count INTEGER,
    last_updated TIMESTAMP DEFAULT CURRENT_TIMESTAMP,
    UNIQUE(ec_number, parameter_type, organism, substrate)
);
\end{verbatim}

\textbf{Design}:

\begin{itemize}
    \item \texttt{brenda\_raw\_data}: Stores every individual measurement (immutable)
    \item \texttt{brenda\_statistics}: Cached aggregated statistics (recomputed on demand)
\end{itemize}

\subsection{Workflow with Caching}
\label{subsec:caching-workflow}

\textbf{Integrated fetcher workflow}:

\begin{enumerate}
    \item Check local DB for cached statistics
    \item If found and fresh (age $<$ 30 days) $\to$ return immediately
    \item Else: Fetch from BRENDA API $\to$ save to DB $\to$ return
\end{enumerate}

\textbf{Example}:

\begin{verbatim}
inferencer = CachedParameterInferencer(brenda_conn, brenda_db, 
                                       organism='Saccharomyces cerevisiae')

# First call: Fetch from BRENDA (2 seconds)
km_glucose = inferencer.infer_km("2.7.1.1", "D-glucose")
# Output: Fetching Km for 2.7.1.1 from BRENDA API...
#         Cached 347 new entries
# Result: {'median': 0.165, 'confidence_interval_95_lower': 0.10, ...}

# Second call: Read from cache (10 ms)
km_glucose_2 = inferencer.infer_km("2.7.1.1", "D-glucose")
# Output: Using cached Km for 2.7.1.1
# Result: {'median': 0.165, ...}  (same, but 200× faster)
\end{verbatim}

\section{UI Integration}
\label{sec:brenda-ui-integration}

\subsection{``Fetch from BRENDA'' Button}
\label{subsec:fetch-from-brenda-button}

\textbf{Transition property panel} includes:

\begin{enumerate}
    \item \textbf{EC Number field}: User enters EC number (e.g., ``2.7.1.1'')
    \item \textbf{``Fetch from BRENDA'' button}: Triggers parameter inference
    \item \textbf{Results table}: Shows median $K_m$, $V_{\max}$, $K_i$ with confidence intervals
    \item \textbf{``Apply'' button}: Fills rate function parameters
\end{enumerate}

\textbf{User workflow}:

\begin{enumerate}
    \item User creates transition (e.g., ``Hexokinase'')
    \item User enters EC number ``2.7.1.1''
    \item User clicks ``Fetch from BRENDA''
    \item SHYpn displays:
    \begin{verbatim}
    Km(glucose):  0.165 mM [0.10-0.22]  (23 values)
    kcat:         100 s⁻¹ [80-120]      (15 values)
    \end{verbatim}
    \item User clicks ``Apply'' $\to$ Rate function auto-configured
\end{enumerate}

\subsection{Batch Parameter Inference}
\label{subsec:batch-parameter-inference}

\textbf{For entire pathways} (e.g., glycolysis with 10 transitions):

\begin{enumerate}
    \item User selects \textbf{``Enrich All Transitions''} menu item
    \item SHYpn iterates over all transitions with EC numbers
    \item For each: Fetch BRENDA parameters $\to$ Update rate function
    \item Progress bar shows ``5/10 transitions enriched''
\end{enumerate}

\section{Limitations and Future Work}
\label{sec:brenda-limitations}

\subsection{Current Limitations}
\label{subsec:brenda-current-limitations}

\begin{enumerate}
    \item \textbf{Missing data}: Not all enzymes in KEGG are in BRENDA
    \begin{itemize}
        \item \textbf{Mitigation}: Fallback to literature search or user input
    \end{itemize}
    
    \item \textbf{$V_{\max}$ vs. $k_{cat}$}: BRENDA provides $k_{cat}$, but $V_{\max} = k_{cat} \cdot [E]_0$ requires enzyme concentration
    \begin{itemize}
        \item \textbf{Solution}: Assume default $[E]_0 = 0.01$ mM (adjustable by user)
    \end{itemize}
    
    \item \textbf{Multi-substrate kinetics}: Most BRENDA data is for single-substrate reactions
    \begin{itemize}
        \item \textbf{Limitation}: Ordered bi-bi kinetics not well-covered
    \end{itemize}
    
    \item \textbf{pH/temperature dependence}: BRENDA values at different conditions
    \begin{itemize}
        \item \textbf{Future}: Temperature correction using Arrhenius equation
    \end{itemize}
\end{enumerate}

\subsection{Future Enhancements}
\label{subsec:brenda-future-enhancements}

\begin{enumerate}
    \item \textbf{Machine learning refinement}:
    \begin{itemize}
        \item Train regression model on cached BRENDA data
        \item Predict $K_m$ for enzymes without BRENDA data (based on substrate structure)
    \end{itemize}
    
    \item \textbf{SABIO-RK integration}:
    \begin{itemize}
        \item Alternative kinetics database with curated models
        \item Often has complete kinetic laws (not just $K_m$)
    \end{itemize}
    
    \item \textbf{Uncertainty propagation}:
    \begin{itemize}
        \item Use confidence intervals in simulation (Monte Carlo sampling)
        \item Report ``parameter uncertainty impacts reachability by $\pm$10\%''
    \end{itemize}
\end{enumerate}

\section{Summary}
\label{sec:brenda-summary}

\textbf{Chapter 10 presented intelligent parameter inference from BRENDA}:

\begin{enumerate}
    \item \textbf{BRENDA SOAP API}: Fetch $K_m$, $k_{cat}$, $K_i$ from 2.7 million measurements
    \item \textbf{Statistical aggregation}: Median + 95\% CI (robust to outliers)
    \item \textbf{Context-aware heuristics}:
    \begin{itemize}
        \item Organism priority (yeast $>$ human $>$ all)
        \item Substrate fuzzy matching
        \item Quality filtering (citation, commentary, outlier detection)
    \end{itemize}
    \item \textbf{Local caching}: SQLite database $\to$ 200$\times$ faster on repeated queries
    \item \textbf{UI integration}: ``Fetch from BRENDA'' button, batch enrichment
    \item \textbf{Performance}: Entire glycolysis pathway (10 enzymes) enriched in $<10$ seconds
\end{enumerate}

\textbf{Key innovation}: \textbf{Progressive data accumulation}

\begin{itemize}
    \item Each BRENDA query adds to local database
    \item Over time, user builds comprehensive organism-specific kinetics library
    \item Enables offline modeling (no internet required after initial fetch)
\end{itemize}

\textbf{Example}: User models yeast metabolism for 1 month $\to$ accumulates 500 EC numbers $\times$ 50 values each = 25,000 cached measurements $\to$ all future yeast models use cached data (instant parameter inference)

\textbf{Next chapter} (Chapter~\ref{ch:simulation}): Hybrid simulation engine (ODE + Gillespie + Timed + Burst).
